\section{Methods}\label{section:methods}

% =====================================================================
%   METHODS / FORMALISM (target: ~3--4 columns; rubric: 20 pts)
%   Required ingredients (per EvaluationForm.md):
%     - Discussion of the methods used and their basis / suitability.
%   Strategy:
%     - Derive the full pipeline from financial data to QAOA circuit.
%     - State every transformation explicitly so the reader can reproduce.
%     - Discuss why each step is a sensible choice and what it costs.
% =====================================================================

The QAOA pipeline maps a financial dataset to a quantum circuit through a sequence of well-defined steps:
\begin{equation}
    \text{Financial data} \rightarrow (\mu,\Sigma) \rightarrow \text{QUBO} \rightarrow \text{Ising} \rightarrow \text{Pauli operators} \rightarrow \text{QAOA}.
\end{equation}
Subsections \ref{subsec:meanvar}--\ref{subsec:qaoa} develop each arrow in turn. Subsection \ref{subsec:xymixer} introduces a constraint-preserving mixer that we compare against the penalty encoding, and Subsection \ref{subsec:metrics} fixes the metrics used in Section~\ref{section:result}.


\subsection{Mean--Variance Portfolio Optimisation}\label{subsec:meanvar}

For a binary selection vector $x \in \{0,1\}^n$ over $n$ candidate assets, the Markowitz objective \autocite{markowitz1952} is
\begin{equation}
    \min_x \left( -\mu^\top x + \lambda\, x^\top \Sigma x \right),
\end{equation}
where $-\mu^\top x$ maximises expected return and $\lambda\, x^\top \Sigma x$ penalises portfolio variance. The risk-aversion parameter $\lambda \geq 0$ tunes the trade-off.

A fixed-budget constraint requires exactly $K$ assets,
\begin{equation}
    \sum_{i=1}^{n} x_i = K,
\end{equation}
which we enforce with a quadratic penalty, giving the soft-constrained cost
\begin{equation}\label{eq:Cx}
    C(x) = -\mu^\top x + \lambda\, x^\top \Sigma x + A \left( \sum_{i=1}^{n} x_i - K \right)^2,
\end{equation}
where $A > 0$ is chosen large enough that the global minimum of $C(x)$ satisfies the budget constraint. Section \ref{subsec:Asensitivity-results} studies the sensitivity of this calibration.

Returns and their first two moments are estimated from price series:
\begin{equation}
    r_i(t) = \frac{P_i(t+1)-P_i(t)}{P_i(t)} \approx \ln\frac{P_i(t+1)}{P_i(t)},
\end{equation}
\begin{equation}
    \mu_i = \langle r_i \rangle_t,
    \qquad
    \Sigma_{ij} = \langle r_i r_j \rangle - \langle r_i \rangle \langle r_j \rangle.
\end{equation}


\subsection{Quadratic Unconstrained Binary Optimisation (QUBO)}\label{subsec:qubo}

A QUBO is a quadratic cost over binary variables,
\begin{equation}
    \min_{x\in\{0,1\}^n} C(x) = \sum_i Q_{ii}\, x_i + \sum_{i<j} Q_{ij}\, x_i x_j,
\end{equation}
where the linear terms are absorbed into the diagonal of $Q$ because $x_i^2 = x_i$ for binary variables. Expanding the penalised mean--variance cost \eqref{eq:Cx} yields
\begin{equation}
    C(x) = \sum_i \underbrace{\bigl[-\mu_i + \lambda\Sigma_{ii} + A(1-2K)\bigr]}_{Q_{ii}}\, x_i + \sum_{i<j} \underbrace{\bigl[2\lambda\Sigma_{ij} + 2A\bigr]}_{Q_{ij}}\, x_i x_j + \underbrace{AK^2}_{\text{const}},
\end{equation}
\begin{equation}
    \boxed{
    Q_{ii} = -\mu_i + \lambda\Sigma_{ii} + A(1-2K),
    \qquad
    Q_{ij} = 2\lambda\Sigma_{ij} + 2A \quad (i<j).
    }
\end{equation}


\subsection{Ising Hamiltonian Mapping}\label{subsec:ising}

Mapping binary variables to spins via
\begin{equation}
    x_i = \frac{1-z_i}{2},\qquad z_i\in\{-1,+1\},
\end{equation}
and substituting into $C(x)$ produces an Ising Hamiltonian
\begin{equation}
    H(z) = \sum_i h_i\, z_i + \sum_{i<j} J_{ij}\, z_i z_j,
\end{equation}
with coefficients
\begin{equation}
    h_i = -\tfrac{1}{2} Q_{ii} - \tfrac{1}{4} \sum_{j\neq i} Q_{ij},
    \qquad
    J_{ij} = \tfrac{1}{4} Q_{ij} \quad (i<j),
\end{equation}
and additive constant
\begin{equation}
    c = AK^2 + \tfrac{1}{2}\sum_i Q_{ii} + \tfrac{1}{4}\sum_{i<j} Q_{ij},
\end{equation}
chosen so that $H(z) + c = C(x)$ for every binary $x$ and its spin image $z$.


\subsection{Quantum Pauli Operators and Cost Hamiltonian}\label{subsec:pauli}

Promoting each classical spin $z_i$ to the single-qubit Pauli-$Z$ operator on site $i$,
\begin{equation}
    z_i \longrightarrow \hat{Z}_i = I^{\otimes i}
    \otimes Z \otimes I^{\otimes (n-i-1)},
\end{equation}
with
\begin{equation}
    Z = \begin{pmatrix} 1 & 0 \\ 0 & -1 \end{pmatrix},
\end{equation}
gives the cost Hamiltonian as the $2^n\times 2^n$ diagonal matrix
\begin{equation}
    \boxed{
    H_C = \sum_i h_i\, \hat{Z}_i + \sum_{i<j} J_{ij}\, \hat{Z}_i \hat{Z}_j.
    }
\end{equation}
Every computational basis state $|z\rangle$ is an eigenstate of $H_C$ with eigenvalue $H(z)$.

The standard \emph{transverse-field} mixing Hamiltonian drives tunnelling between bitstring states,
\begin{equation}\label{eq:HM-transverse}
    H_M^{(X)} = \sum_{i=1}^{n} \hat{X}_i,
    \qquad
    X = \begin{pmatrix} 0 & 1 \\ 1 & 0 \end{pmatrix}.
\end{equation}
Two properties make $H_M^{(X)}$ a natural choice. First, $\hat X_i$ flips the value of qubit $i$ in the computational basis, so $[H_M^{(X)}, H_C]\neq 0$ and the mixer generates the non-diagonal dynamics that QAOA needs to interfere between candidate bitstrings. Second, the (degenerate) ground state of $-H_M^{(X)}$ is the uniform superposition
\begin{equation}\label{eq:plus-state}
    |s\rangle \;=\; |+\rangle^{\otimes n} \;=\; \frac{1}{\sqrt{2^n}}\sum_{x\in\{0,1\}^n} |x\rangle,
\end{equation}
which is prepared by a single layer of Hadamard gates acting on $|0\rangle^{\otimes n}$ and serves as the initial state of the variational circuit defined below.


\subsection{QAOA Algorithm}\label{subsec:qaoa}

\paragraph{Adiabatic intuition.}
QAOA \autocite{farhi2014qaoa} is most naturally motivated as a discretised, variational analogue of quantum adiabatic computation. The adiabatic theorem states that a system prepared in the ground state of an initial Hamiltonian and slowly evolved towards a final Hamiltonian remains in the instantaneous ground state throughout the evolution. Choosing the final Hamiltonian to be the diagonal cost operator $H_C$ of Section~\ref{subsec:pauli} and the initial Hamiltonian to be the transverse-field mixer $H_M^{(X)}$ — whose ground state is the easy-to-prepare uniform superposition \eqref{eq:plus-state} — converts combinatorial optimisation into ground-state preparation. The total adiabatic evolution time can in the worst case scale exponentially in $n$, however, so QAOA instead approximates the evolution by a fixed number $p$ of short alternating steps generated by $H_C$ and $H_M^{(X)}$ and treats the step durations $(\gamma_k, \beta_k)$ as variational parameters. In the limit $p \to \infty$ with an appropriate Trotterised schedule the ansatz recovers the exact adiabatic algorithm.

The single-layer ($p = 1$) ansatz is
\begin{equation}
    |\psi(\gamma,\beta)\rangle =
    e^{-i \beta H_M}\,
    e^{-i \gamma H_C}\,
    |+\rangle^{\otimes n},
\end{equation}
where $e^{-i\gamma H_C}$ is the \emph{phase-separation} unitary --- it phases each basis state by $e^{-i\gamma H(z)}$, encoding cost information into the amplitudes --- and $e^{-i\beta H_M}$ is the \emph{mixing} unitary, which redistributes amplitude between bitstrings to enable transitions between candidate solutions.

\begin{equation}
    \underbrace{H^{\otimes n}}_{\text{superposition}}
    \longrightarrow
    \underbrace{e^{-i\gamma H_C}}_{\text{phase separation}}
    \longrightarrow
    \underbrace{e^{-i\beta H_M}}_{\text{mixing}}
    \longrightarrow
    \underbrace{\text{measure}}_{Z\text{-basis}}
\end{equation}

For depth $p > 1$, $p$ cost and $p$ mixer layers are alternated,
\begin{equation}\label{eq:multilayer}
    |\psi(\boldsymbol\gamma,\boldsymbol\beta)\rangle =
    \prod_{k=1}^{p}
    \left(
        e^{-i\beta_k H_M}\,
        e^{-i\gamma_k H_C}
    \right)
    |+\rangle^{\otimes n},
\end{equation}
with $2p$ variational parameters $(\gamma_1,\dots,\gamma_p,\beta_1,\dots,\beta_p)$.

\paragraph{Interference mechanism.}
The two unitaries play complementary roles. The diagonal $e^{-i\gamma H_C}$ phases each basis state by $e^{-i\gamma H(z)}$, imprinting the cost landscape onto the wavefunction without changing measurement probabilities. The non-diagonal $e^{-i\beta H_M}$ mixes amplitude between basis states; because $[H_C, H_M] \neq 0$, the two operations together produce non-trivial quantum interference. With the right angles $(\boldsymbol\gamma, \boldsymbol\beta)$, amplitude accumulates constructively on low-cost bitstrings and destructively on high-cost ones, concentrating measurement probability on near-optimal solutions. This is the mechanism by which QAOA differs from a purely classical heuristic such as simulated annealing or tabu search: amplitude flows non-locally across the cost landscape in superposition, whereas classical samplers visit one configuration at a time.

\paragraph{Hybrid quantum--classical loop.}
A classical optimiser minimises the energy expectation value
\begin{equation}
    E(\boldsymbol\gamma, \boldsymbol\beta) =
    \langle \psi(\boldsymbol\gamma,\boldsymbol\beta) |\, H_C\, | \psi(\boldsymbol\gamma,\boldsymbol\beta)\rangle.
\end{equation}
Since $H_C$ is diagonal in the computational basis with eigenvalues $H(z)$, this reduces to
\begin{equation}\label{eq:E-as-prob}
    E(\boldsymbol\gamma, \boldsymbol\beta) = \sum_{z\in\{-1,+1\}^n} H(z)\, |\langle z | \psi \rangle|^2,
\end{equation}
so the QAOA cost is a probability-weighted sum of Ising energies. The quantum circuit has only $2p$ parameters and depth $O(p)$, which is what makes QAOA a prototypical NISQ algorithm: the quantum device contributes a non-classical sampling primitive, while the parameter search itself sits on a classical CPU. The variational principle guarantees $E(\boldsymbol\gamma,\boldsymbol\beta) \geq E_0$ for all angles, with equality recovered in the limit $p \to \infty$ \autocite{farhi2014qaoa}; finite-$p$ QAOA returns an approximation whose quality we track with the metrics of Section~\ref{subsec:metrics}. In this work the classical optimiser is COBYLA with random restarts; the optimisation pipeline is described in Section~\ref{subsec:codestructure}.


\subsection{XY-Mixer: Constraint-Preserving Encoding}\label{subsec:xymixer}

% --- Motivation -------------------------------------------------------
The penalty term $A(\sum_i x_i - K)^2$ in \eqref{eq:Cx} enforces the budget only in the limit $A \to \infty$; for finite $A$, infeasible bitstrings with $\sum_i x_i \neq K$ remain in the search space and may carry non-negligible amplitude after a finite-depth circuit. An alternative is to discard the penalty and instead choose a mixer that preserves Hamming weight by construction \autocite{mancilla2026constrained}. The \emph{XY-ring mixer} is
\begin{equation}\label{eq:HM-XY}
    H_M^{(XY)} = \tfrac{1}{2} \sum_{i=1}^{n} \bigl( \hat{X}_i \hat{X}_{i+1} + \hat{Y}_i \hat{Y}_{i+1} \bigr),
\end{equation}
with periodic boundary conditions ($i + 1$ taken mod $n$) and
\begin{equation}
    Y = \begin{pmatrix} 0 & -i \\ i & 0 \end{pmatrix}.
\end{equation}
The operator $\hat{X}_i \hat{X}_{i+1} + \hat{Y}_i \hat{Y}_{i+1}$ swaps the basis states $|01\rangle$ and $|10\rangle$ on neighbouring qubits while leaving $|00\rangle$ and $|11\rangle$ invariant, so it commutes with the Hamming-weight operator
\begin{equation}
    \hat{N} = \sum_i \tfrac{1}{2}(I - \hat{Z}_i).
\end{equation}
When the initial state is a feasible Dicke state $|D^n_K\rangle$ --- the equal superposition over all bitstrings with $\sum_i x_i = K$ --- the entire QAOA trajectory stays in the feasible subspace and the penalty parameter $A$ becomes unnecessary.

% --- Trade-offs to discuss in the Results section --------------------
We compare the two encodings on the same data in Section~\ref{subsec:xymixer-results}: the penalty mixer is simpler to implement but inherits a hyperparameter; the XY-mixer is hyperparameter-free but requires Dicke-state preparation and has a slightly larger gate count per layer.


\subsection{Statistical Physics Connection}\label{subsec:statphys}

The Ising form of $H_C$ defines a classical partition function at inverse temperature $\beta_{\text{th}}$,
\begin{equation}
    Z = \sum_{\{z\}} e^{-\beta_{\text{th}} H(z)},
\end{equation}
with associated free energy
\begin{equation}
    F = -\frac{1}{\beta_{\text{th}}} \ln Z.
\end{equation}
QAOA can be read as a variational approximation to the ground state of $H_C$,
\begin{equation}
    E(\boldsymbol\gamma, \boldsymbol\beta) = \langle \psi | H_C | \psi \rangle \geq E_0,
\end{equation}
and as $p \to \infty$ the ansatz can in principle reach the exact ground state. Finite-$p$ QAOA is thus directly analogous to a finite-temperature thermal state of the same Ising Hamiltonian, a connection we explore numerically in Section~\ref{subsec:thermal-results}.


\subsection{Metrics and Validation}\label{subsec:metrics}

To compare runs across datasets and circuit depths we use:

\paragraph{Approximation ratio.}
Following the QAOA literature \autocite{farhi2014qaoa,zhou2018qaoa}, we define
\begin{equation}\label{eq:approxratio}
    r(p) = \frac{E(\boldsymbol\gamma^\star, \boldsymbol\beta^\star)\big|_{p} - C_{\max}}{C_{\min} - C_{\max}} \in [0, 1],
\end{equation}
where $C_{\min} = \min_x C(x)$ and $C_{\max} = \max_x C(x)$ are obtained by exhaustive enumeration at $n = 4$. A value $r = 1$ corresponds to recovering the optimum in expectation.

\paragraph{Top-1 success probability.}
The probability that a single shot returns the true optimum,
\begin{equation}
    p_{\text{top-1}}(p) = |\langle z^\star | \psi(\boldsymbol\gamma^\star, \boldsymbol\beta^\star) \rangle|^2,
\end{equation}
with $z^\star$ the brute-force optimum.

\paragraph{Feasibility ratio.}
The probability that a measurement yields a budget-feasible bitstring,
\begin{equation}
    p_{\text{feas}}(p) = \sum_{z : \sum_i x_i = K} |\langle z | \psi \rangle|^2,
\end{equation}
which equals $1$ identically under the XY-mixer and is a non-trivial diagnostic under the penalty mixer.

\paragraph{Validation strategy.}
At $n = 4$ the brute-force optimum is computable in $2^4 = 16$ evaluations, so every QAOA run is anchored against the exact answer. Our QAOA implementation is additionally cross-checked against the construction in the FYS5419 course material at $p = 1$.
